% GENERATED FILE. Edit canonical lessons, then run npm run generate:books.
% canonical-source-hash: fnv1a64:e9e99b5b099b6c02
% canonical-lessons: PT-C441-campanario, PT-C441-catedral, PT-C441-arcadas, PT-C441-chafariz, PT-C441-pelourinho

\chapter{Na praça}
\label{ch:pt-na-praca}
\hlchaptermodality{441}

\begin{chapteropening}
\textbf{By the end of this chapter:} \emph{I can describe the buildings round a square: the bell tower, the cathedral, the arcades, the fountain and the stone pillory.}

Complete the last lesson of chapter 441: I can describe the buildings round a square: the bell tower, the cathedral, the arcades, the fountain and the stone pillory.
\end{chapteropening}

\section[o campanário]{o campanário — the bell tower}

\label{lesson:PT-C441-campanario}



\begin{quote}
Before the new one: say the Portuguese for the municipality, then the Portuguese for the entry phone.
\end{quote}

\subsection*{You'll want to know: o campanário}
\textbf{o campanário} — \textquotedblleft{}the bell tower\textquotedblright{}.

\begin{quote}
Do campanário vê-se toda a cidade.
\end{quote}

\begin{quote}
\emph{From the bell tower you can see all of the town.}
\end{quote}

Built on Latin \emph{campāna}, a bell.

\begin{cousinweb}
\emph{campanário} is a learned formation on Late Latin \emph{campāna} 'bell' (Lewis \& Short).
\end{cousinweb}

\begin{grammarlens}[title={What you've built}]
One more word for telling, explaining and describing at B1.
\end{grammarlens}

\subsection*{Your turn}
\begin{itemize}
\raggedright
  \item \emph{Say it:} \emph{o campanário}
  \item \emph{Say it:} \emph{o campanário}, once more
  \item \emph{Say it:} a sentence of your own with \emph{o campanário}
  \item \emph{Recall:} say the Portuguese for the municipality, then the Portuguese for the entry phone, then say \emph{o campanário} again
\end{itemize}

\begin{tcolorbox}[breakable,colback=teal!4,colframe=teal!35!black,title={Before you move on}]
What does o campanário mean? (\textquotedblleft{}The bell tower\textquotedblright{}.) Say it once more.
\end{tcolorbox}
\section[a catedral]{a catedral — the cathedral}

\label{lesson:PT-C441-catedral}



\begin{quote}
Before the new one: say the Portuguese for the entry phone, then the Portuguese for the bell tower.
\end{quote}

\subsection*{You'll want to know: a catedral}
\textbf{a catedral} — \textquotedblleft{}the cathedral\textquotedblright{}.

\begin{quote}
A catedral é o edifício mais antigo da praça.
\end{quote}

\begin{quote}
\emph{The cathedral is the oldest building on the square.}
\end{quote}

In Portugal a cathedral is often called \emph{a Sé}.

\begin{cousinweb}
\emph{catedral} is a learned borrowing from Late Latin \emph{cathedrālis} 'of the bishop's seat', built on \emph{cathedra} 'chair', from Greek \emph{kathédra} 'seat' (Lewis \& Short has \emph{cathedra}).
\end{cousinweb}

\begin{grammarlens}[title={What you've built}]
One more word for telling, explaining and describing at B1.
\end{grammarlens}

\subsection*{Your turn}
\begin{itemize}
\raggedright
  \item \emph{Say it:} \emph{a catedral}
  \item \emph{Say it:} \emph{a catedral}, once more
  \item \emph{Say it:} a sentence of your own with \emph{a catedral}
  \item \emph{Recall:} say the Portuguese for the entry phone, then the Portuguese for the bell tower, then say \emph{a catedral} again
\end{itemize}

\begin{tcolorbox}[breakable,colback=teal!4,colframe=teal!35!black,title={Before you move on}]
What does a catedral mean? (\textquotedblleft{}The cathedral\textquotedblright{}.) Say it once more.
\end{tcolorbox}
\section[as arcadas]{as arcadas — the arcades}

\label{lesson:PT-C441-arcadas}



\begin{quote}
Before the new one: say the Portuguese for the bell tower, then the Portuguese for the cathedral.
\end{quote}

\subsection*{You'll want to know: as arcadas}
\textbf{as arcadas} — \textquotedblleft{}the arcades\textquotedblright{}.

\begin{quote}
Debaixo das arcadas há cafés e lojas.
\end{quote}

\begin{quote}
\emph{Under the arcades there are cafés and shops.}
\end{quote}

Singular \emph{a arcada} — a single arch or arcade.

\begin{cousinweb}
\emph{arcada} is built on \emph{arco} 'arch', which is inherited from Latin \emph{arcus} 'bow, arch' (EtymDB).
\end{cousinweb}

\begin{grammarlens}[title={What you've built}]
One more word for telling, explaining and describing at B1.
\end{grammarlens}

\subsection*{Your turn}
\begin{itemize}
\raggedright
  \item \emph{Say it:} \emph{as arcadas}
  \item \emph{Say it:} \emph{as arcadas}, once more
  \item \emph{Say it:} a sentence of your own with \emph{as arcadas}
  \item \emph{Recall:} say the Portuguese for the bell tower, then the Portuguese for the cathedral, then say \emph{as arcadas} again
\end{itemize}

\begin{tcolorbox}[breakable,colback=teal!4,colframe=teal!35!black,title={Before you move on}]
What does as arcadas mean? (\textquotedblleft{}The arcades\textquotedblright{}.) Say it once more.
\end{tcolorbox}
\section[o chafariz]{o chafariz — the public fountain}

\label{lesson:PT-C441-chafariz}



\begin{quote}
Before the new one: say the Portuguese for the cathedral, then the Portuguese for the arcades.
\end{quote}

\subsection*{You'll want to know: o chafariz}
\textbf{o chafariz} — \textquotedblleft{}the public fountain\textquotedblright{}.

\begin{quote}
Antigamente, as pessoas iam buscar água ao chafariz.
\end{quote}

\begin{quote}
\emph{In the old days, people went to fetch water from the fountain.}
\end{quote}

A stone fountain where people once filled their jugs.

\begin{cousinweb}
\emph{chafariz} comes from Arabic \emph{ṣihrīj} 'cistern, water tank' (EtymDB), like Catalan \emph{safareig}.
\end{cousinweb}

\begin{grammarlens}[title={What you've built}]
One more word for telling, explaining and describing at B1.
\end{grammarlens}

\subsection*{Your turn}
\begin{itemize}
\raggedright
  \item \emph{Say it:} \emph{o chafariz}
  \item \emph{Say it:} \emph{o chafariz}, once more
  \item \emph{Say it:} a sentence of your own with \emph{o chafariz}
  \item \emph{Recall:} say the Portuguese for the cathedral, then the Portuguese for the arcades, then say \emph{o chafariz} again
\end{itemize}

\begin{tcolorbox}[breakable,colback=teal!4,colframe=teal!35!black,title={Before you move on}]
What does o chafariz mean? (\textquotedblleft{}The public fountain\textquotedblright{}.) Say it once more.
\end{tcolorbox}
\section[o pelourinho]{o pelourinho — the pillory (stone pillar)}

\label{lesson:PT-C441-pelourinho}



\begin{quote}
Before the new one: say the Portuguese for the arcades, then the Portuguese for the public fountain.
\end{quote}

\subsection*{You'll want to know: o pelourinho}
\textbf{o pelourinho} — \textquotedblleft{}the pillory (stone pillar)\textquotedblright{}.

\begin{quote}
No centro da praça há um pelourinho de pedra.
\end{quote}

\begin{quote}
\emph{In the middle of the square there is a stone pillory.}
\end{quote}

A carved stone pillar that showed a town had its own courts.

\begin{cousinweb}
\emph{pelourinho} comes from French \emph{pilori} 'pillory' (EtymDB), reshaped with the small-size ending \emph{-inho}; the French word's origin is not settled.
\end{cousinweb}

\begin{grammarlens}[title={What you've built}]
One more word for telling, explaining and describing at B1.
\end{grammarlens}

\subsection*{Your turn}
\begin{itemize}
\raggedright
  \item \emph{Say it:} \emph{o pelourinho}
  \item \emph{Say it:} \emph{o pelourinho}, once more
  \item \emph{Say it:} a sentence of your own with \emph{o pelourinho}
  \item \emph{Recall:} say the Portuguese for the arcades, then the Portuguese for the public fountain, then say \emph{o pelourinho} again
\end{itemize}

\begin{tcolorbox}[breakable,colback=teal!4,colframe=teal!35!black,title={Before you move on}]
What does o pelourinho mean? (\textquotedblleft{}The pillory (stone pillar)\textquotedblright{}.) Say it once more.
\end{tcolorbox}
